%include directives.fmt

% Some freest specific formatting directives
%format linS = "\mathrm{1S}"

\appendix


The organization of the supplement follows the structure of the main document.
Topics are covered in the same sequence, so the two documents can be read
side-by-side.


\section{Motivation --- Remarks}
\label{sec:remarks}

%\subsection{Remarks}

% Switch to the borrow code style for this appendix section
{\BorrowCodeStyle

The reader may ask why |select| still adheres to
resource-passing style. The answer is simplicity of our system. While it
is possible to define |select| with a borrowing-based
interface, supporting it requires subtyping of session types. While
subtyping is fine in theory, it turns out that it is undecidable for
CFST~\cite{DBLP:journals/toplas/Padovani19}. Hence, it cannot be the
basis for a complete type checker.

For borrowing from regular session
types~\cite{DBLP:journals/pacmpl/SaffrichS0V25}, a borrowing-style
|select| operation is possible because subtyping for regular
session types is decidable.

A borrowing-style |branch| operator (external choice) cannot be supported in \thecalculus.
This operator selects between different continuations at run time and
the residual channel depends on this run-time choice.

% The downside is that some uses of |select| force us to use
% the |drop| primitive, which discharges an exhausted channel
% borrow of type |Drop|. The discharge is part of the
% |send| and |recv| operations, so it is rarely
% needed in practice. Alternatively, the type checker could insert
% |drop| automatically.

All \thecalculus examples are available as part of our supplement. Compared to the
presentation in the main text, the split operators carry type
parameters in the actual code to simplify type checking.

% End the borrow code style for this appendix section.
}

\clearpage

\section{Types: Session-Type Laws and Well Formedness Predicates}
\label{sec:appendix-types}

\begin{figure}[h]
  \begin{mathpar}
    \RuleUUnit \and
    \RuleUArr \and
    \RuleUProd \and
    \RuleUVariant \\
    \RuleUCtxEmpty \and
    \RuleUCtxVar \and
    \RuleUCtxSeq \and
    \RuleUCtxPar \\
    \RuleSNewSkip \and
    \RuleSNewSend \and
    \RuleSNewRecv \and
    \RuleSNewVar \\
    \RuleSNewSeq \and
    \RuleSNewBranch \and
    \RuleSNewChoice \and
    \RuleSNewRec \\
    \RuleSLocalAcq \and
    \RuleSLocalSeqI \and
    \RuleSLocalSeqII \and
    \RuleSLocalRec \\
    {\SLocal\S \Longleftrightarrow \SNonLocal\S \text{~not derivable}}
  \end{mathpar}
  \caption{Unrestricted types and contexts ($\Unr\T$, $\Unr\Ctx$); session
    types valid for $\CNew$ ($\SNew\S$); local and non-local session types
    ($\SLocal\S$, $\SNonLocal\S$)}
  \label{fig:omitted-predicates}
\end{figure}

\vfill

\begin{figure}[h]
  \medskip
  Laws for context-free session types.
  \begin{mathpar}
    \SSeq\SSkip\S = \S \and
    \SSeq \S\SSkip = \S \and
    \SSeq{\S[1]}{(\SSeq{\S[2]}{\S[3]})} =
    \SSeq{(\SSeq{\S[1]}{\S[2]})}{{\S[3]}} \and
    \SRec\SVar\S = \S{[\SRec\SVar\S/\SVar]} \and
    \SSeq{(\SChoice{\ell:\S[\ell]}{\ell\in L})}\S =
    \SChoice{\ell:\SSeq{\S[\ell]}\S}{\ell\in L} \and
    \SSeq{(\SBranch{\ell:\S[\ell]}{\ell\in L})}\S = \SBranch{\ell:\SSeq{\S[\ell]}\S}{\ell\in L}
  \end{mathpar}

  \medskip
  Laws for duality $\SDual\S$ (there is no dual for $\SDrop$ and $\SAcq$).
  \begin{mathpar}
    \SDual\SSkip = \SSkip \and
    \SDual{(\SSeq{\S[1]}{\S[2]})} = \SSeq{\SDual{\S[1]}}{\SDual{\S[2]}}
    \and
    \SDual{\SSend\T} = \SRecv\T \and
    \SDual{\SRecv\T} = \SSend\T \and
    \SDual{(\SChoice{\ell:\S[\ell]}{\ell\in L})} =
    \SBranch{\ell:\SDual{\S[\ell]}}{\ell\in L} \and
    \SDual{(\SBranch{\ell:\S[\ell]}{\ell\in L})} =
    \SChoice{\ell:\SDual{\S[\ell]}}{\ell\in L} \and
    \SDual\STerm = \SWait \and
    \SDual\SWait = \STerm \and
    \SDual{(\SRec\SVar\S)} = \SDual{\S{[\SRec\SVar\S/\SVar]}}
  \end{mathpar}
  \caption{Session-type equality and duality}
  \label{fig:session-laws}
\end{figure}

\vfill
\clearpage

\section{Dynamics: Congruence and Reduction Rules}

\begin{figure}[h]
\begin{align*}
  \U,\V \grmeq
    & \C
      \grmor \EVar
      \grmor \EAbs \EVar\E
      \grmor \EInj{i} \V
      \grmor \EPair \U\V
      % \grmor \CSend \V 
              \tag{Values}
  \\
  \EC \grmeq&
              \ECEmpty
              \grmor \EAppD[\DirLeft] \E\EC
              \grmor \EAppD[\DirLeft] \EC\V
              \grmor \EAppD[\DirRight] \EC\E
              \grmor \EAppD[\DirRight] \V\EC
              \grmor \ESeq \EC\E
              \grmor \EPair \EC\E
              \grmor \EPair \V\EC
              \\ \grmor
    & \EInj{i}\EC \grmor \ELet {\EPair \EVar\EVar}\EC\E
              \grmor \ECase \EC {\ECaseBranch* i \EVar \E}
              \tag{Evaluation contexts}
  \\
  \FC \grmeq
  & \PExp \EC
    \tag{Process evaluation contexts}
  \\
  \CC \grmeq
  & \ECEmpty
    \grmor \PPar{\CC}{\Proc}
    \grmor \PNew \BindGroup\BindGroup \CC
    \tag{Process contexts}
\end{align*}
  \caption{Values and evaluation contexts}
  \label{fig:values-evaluation-contexts}
\end{figure}

\vfill

\begin{figure}[h]
  \begin{mathpar}
    \RuleExpRedApp \and
    \RuleExpRedSeq \and
    \RuleExpRedPairElim \and
    \RuleExpRedUnfold \and
    \RuleExpRedVariantElim \and
    \RuleExpRedCtx
  \end{mathpar}
  \caption{Expression reduction ($\E \EReduce \E$)}
  \label{fig:expression-reduction}
\end{figure}

\vfill

\begin{figure}[h]
  \begin{mathpar}
    % \RuleProcCongRefl \and
    % \RuleProcCongTrans \and
    \RuleProcCongParComm \and
    \RuleProcCongParAssoc \and
    \RuleProcCongParUnit \and
    \RuleProcCongResSwap \and
    \RuleProcCongResComm \and
    \RuleProcCongExtend
  \end{mathpar}
  \caption{Process congruence}
  \label{fig:configuration-congruence}
\end{figure}

\vfill
\clearpage


\section{Algorithmic Typing}
\label{sec:algorithmic-typing-appendix}

\begin{figure}[h]
  \begin{minipage}{0.49\linewidth}
  \[
    \Ctx[1] \CtxJoin \Ctx[2] = \begin{cases}
      \Ctx[1] \CtxPar \Ctx[2] & \text{if $\Dir \in \{ \DirUnord, \DirNone \}$} \\
      \Ctx[1] \CtxSeq \Ctx[2] & \text{if $\Dir = \DirLeft$} \\
      \Ctx[2] \CtxSeq \Ctx[1] & \text{if $\Dir = \DirRight$}
    \end{cases}
  \]
  \end{minipage}
  \begin{minipage}{0.49\linewidth}
  \[
    \CtxRestrict*{X} = \begin{cases}
      \CtxRestrict*[1]{X} \CtxPar \CtxRestrict*[2]{X} & \text{if $\Ctx=\Ctx[1] \CtxPar \Ctx[2]$} \\
      \CtxRestrict*[1]{X} \CtxSeq \CtxRestrict*[2]{X} & \text{if $\Ctx=\Ctx[1] \CtxSeq \Ctx[2]$} \\
      \CtxVar \EVar \T & \text{if $\Ctx=\CtxVar \EVar T$ and $\EVar \in X$} \\
      \CtxEmpty & \text{otherwise}
    \end{cases}
  \]
  \end{minipage}
  \caption{Context join $\Ctx[1] \CtxJoin \Ctx[2]$ joins two contexts
    according to the direction $\Dir$, and context restriction
    $\CtxRestrict*{X}$ restricts a context to a set of variables $X$.}
  \label{fig:context-operators}
\end{figure}

\begin{figure}[h]
  \setlength{\belowdisplayskip}{0pt}
  \small
  \hfill \fbox{$\InferType \E \T \Eff$}
  \begin{mathpar}
    \RuleAlgVar \and
    \RuleAlgConst \and
    \RuleAlgAnnot \and
    \RuleAlgLSplit \and
    \RuleAlgRSplit \and
    \RuleAlgApp \and
    \RuleAlgSeq \and
    \RuleAlgLetPair \and
    \RuleAlgCase
  \end{mathpar}
  \vspace{\abovecaptionskip}
  \[ \DefinitionsReqDirEff \]
  \caption{Algorithmic typing, inference mode}
  \label{fig:appendix-algorithmic-infer}
\end{figure}

\begin{figure}[hp]
  \small
  \vspace{1em}
  \hfill \fbox{$\CheckType \E \T \Eff$}
  \begin{mathpar}
    \RuleAlgAbs \and
    \RuleAlgCheck \and
    \RuleAlgAbsRec \and
    \RuleAlgPair \and
    \RuleAlgInj 
  \end{mathpar}
  \caption{Algorithmic typing, checking mode}
  \label{fig:appendix-algorithmic-check}
\end{figure}

\begin{figure}[h]
  \begin{mathpar}
    \RuleCSPair \and
    \RuleCSSum \and
    \RuleCSUnit \and
    \RuleCSArrow \and
    \RuleCSSwap \and
    \RuleCSContext
  \end{mathpar}
  \caption{Constraint simplification}
  \label{fig:constraint-simplification}
\end{figure}


\clearpage
\section{Implementation}
\label{sec:appendix-implementation}

\subsection{Constraint Checking with FreeST}

In the implementation, we use FreeST
\cite{DBLP:journals/corr/abs-1904-01284,DBLP:journals/iandc/AlmeidaMTV22,
  almeida26:_frees_concur_progr_languag_with} version 5.0 to perform equivalence
checking. To this end, we generate little FreeST programs that typecheck if and
only if the equivalence constraints hold. 


For every equivalence constraint $\T[1] \CEq \T[2]$, we generate the following
FreeST program consisting of type definitions for $\T[1]$ and $\T[2]$ as well as
an identity function mapping $\T[1]$ to $\T[2]$.
For example, let $\T[1] = \SRec{\SVar}{\SSeq {\SSeq {\SSend{\BInt}} {\SSend {\BString}}} {\SVar}}$ and $\T[2] = \SRec{\SVar}{\SSeq {\SSkip} {\SSeq {\SSeq {\SSend{\BInt}} {\SSend {\BString}}} {\SVar}}}$. For the constraint $\T[1] \CEq \T[2]$ we generate the following FreeST program:

{
  \setlength{\blanklineskip}{\medskipamount}
  \TightCode[\bigskipamount]
\begin{code}
^^^^  type X_T1 : linS
      type X_T1 = !Int;!String;X_T1

      type T1 : linS
      type T1 = X_T1

      type X_T2 : linS
      type X_T2 = Skip;!Int;!String;X_T2

      type T2 : linS
      type T2 = X_T2

      id : T1 -> T2
      id x = x
\end{code}
}

\noindent FreeST 5.0 requires specifying kinds for type definitions.
Kind |linS| says that the type needs to be a session type.
If |T1| and |T2| are the translated types from the equivalence constraint,
it is easy to see that this program only typechecks, if |T1| and
|T2| are equivalent session types.

The main parts of the translation are:
\begin{itemize}
  \item Translating inline recursion of the form $\SRec\SVar\S$ to explicit recursion.
  \item Translating mobility and direction annotations on function and pair
    types.
  \item Translating effects.
\end{itemize}

To translate a $\mu$-type to explicit recursion, we generate a fresh name and
associate the name with the recursion unrolled and translation applied.
We return the associated name, instead of the type.
We encode mobility, direction, and effects in a branch type with labels
indicating the state of the annotation.
Thus, the translation function $\tofreest{\T}$, defined in
\cref{fig:freest-translation}, takes a type $\T$ and returns a pair
$\langle F,U\rangle$ consisting of a FreeST type $F$ and a list of explicit
type definitions $U$.

\begin{figure}[tp]
  \footnotesize
  \setlength{\jot}{0.1ex}
  \begin{align*}
    \tofreest{\TypeKeyword{Unit}}
      & = \langle (), \emptyset \rangle \\
    \tofreest{\SSkip}
      & = \langle \SSkip, [] \rangle \\
    \tofreest{\STerm}
      & = \langle \STerm, [] \rangle \\
    \tofreest{\SWait}
      & = \langle \SWait, [] \rangle \\
    \tofreest{\TArr{\Mob}{\Dir, {\Eff}}{\T[1]}{\T[2]}}
      & = \langle (F_1 \to F_2) \TPair \SBranch {\ell : \SSkip} {\ell \in \text{dir} \cup \text{mob} \cup \text{eff}},\,
        U_1 \cup U_2 \rangle \quad \text{where} \\
      & \quad \langle F_1, U_1 \rangle = \tofreest{\T[1]} \\
      & \quad \langle F_2, U_2 \rangle = \tofreest{\T[2]} \\
      & \text{dir} = \begin{cases}
        \{ \text{left} \} & \text{if } \Dir = \DirLeft \\
        \{ \text{right} \} & \text{if } \Dir = \DirRight \\
        \{ \text{lin} \} & \text{if } \Dir = \DirUnord \\
        \emptyset & \text{if } \Dir = \DirNone \\
      \end{cases} \\
      & \text{mob} = \begin{cases}
        \{ \text{mobile} \} & \text{if } \Mob = \MobMobile \\
        \{ \text{static} \} & \text{if } \Mob = \MobStationary \\
      \end{cases} \\
      & \text{eff} = \begin{cases}
        \{ \text{impure} \} & \text{if } \Eff = \EffImpure \\
        \{ \text{} \} & \text{if } \Eff = \EffImpure \\
      \end{cases} \\
    \tofreest{\T[1] \TPair[\Dir] \T[2]}
      & = \langle (F_1 \times F_2) \TPair \SBranch {\ell : \SSkip} {\ell \in \text{dir}},\,
        U_1 \cup U_2 \rangle \quad \text{where} \\
      & \quad \langle F_1, U_1 \rangle = \tofreest{\T[1]} \\
      & \quad \langle F_2, U_2 \rangle = \tofreest{\T[2]} \\
      & \text{dir} = \begin{cases}
        \{ \text{left} \} & \text{if } \Dir = \DirLeft \\
        \{ \text{right} \} & \text{if } \Dir = \DirRight \\
        \{ \text{lin} \} & \text{if } \Dir = \DirUnord \\
        \emptyset & \text{if } \Dir = \DirNone
      \end{cases} \\
    \tofreest{\TVariant{\ell:\T[\ell]}}
      & = \langle \TVariant{\ell:F_\ell},\, \bigcup_{\ell\in L} U_\ell \rangle \quad \text{where} \\
      & \quad \langle F_\ell, U_\ell \rangle = \tofreest{\T[\ell]} \quad \text{for every } \ell \in L \\
    \tofreest{\SSend\T}
      & = \langle \SSend{F}, U \rangle \quad \text{where} \\
      & \quad \langle F, U \rangle = \tofreest{\T} \\
    \tofreest{\SRecv\T}
      & = \langle \SRecv{F}, U \rangle \quad \text{where} \\
      & \quad \langle F, U \rangle = \tofreest{\T} \\
    \tofreest{\S[1] ; \S[2]}
      & = \langle F_1 ; F_2,\, U_1 \cup U_2 \rangle \quad \text{where} \\
      & \quad \langle F_1, U_1 \rangle = \tofreest{\S[1]} \\
      & \quad \langle F_2, U_2 \rangle = \tofreest{\S[2]} \\
    \tofreest{\SChoice{\ell:\S[\ell]}{\ell\in L}}
      & = \langle +\{\ell : F_\ell\}_{\ell\in L},\, \bigcup_{\ell\in L} U_\ell \rangle \quad \text{where} \\
      & \quad \langle F_\ell, U_\ell \rangle = \tofreest{\S[\ell]} \quad \text{for every } \ell \in L \\
    \tofreest{\SBranch{\ell:\S[\ell]}{\ell\in L}}
      & = \langle \&\{\ell : F_\ell\}_{\ell\in L},\, \bigcup_{\ell\in L} U_\ell \rangle \quad \text{where} \\
      & \quad \langle F_\ell, U_\ell \rangle = \tofreest{\S[\ell]} \quad \text{for every } \ell \in L \\
    \tofreest{\SDrop}
      & = \langle \SSend{\TypeKeyword{Drop}},\, \emptyset \rangle \\
    \tofreest{\SAcq}
      & = \langle \SRecv{\TypeKeyword{Drop}}, \emptyset \rangle \\
    \tofreest{\SRec\SVar\S}
      & = \langle \SVar',\, U \cup \{\SVar' \mapsto U\} \rangle \quad \text{where } \SVar' \text{ is fresh and} \\
      & \quad \langle F, U \rangle = \tofreest{\S\,[\SVar'/\SVar]} \\
    \tofreest{\SVar}
      & = \langle \SVar, \emptyset \rangle
  \end{align*}
  \caption{Translation of types to FreeST types.}
  \label{fig:freest-translation}
\end{figure}

\subsection{Soundness and the Completeness of the translation}

\begin{definition}
  The equivalence relation $\simeq$ on \thecalculus types is the least
  congruence closed under the equivalence rules from
  \citet{DBLP:conf/icfp/ThiemannV16} extended with rules from
  \Cref{fig:cfstb-equivalence}. 
\end{definition}

\begin{figure}[tp]
  \begin{mathpar}
    \inferrule{
      \T[1] \simeq \T*[1] \\
      \T[2] \simeq \T*[2]
    }{
      \TArr{\Mob}{\Dir,\Eff}{\T[1]}{\T[2]}
      \simeq
      \TArr{\Mob}{\Dir,\Eff}{\T*[1]}{\T*[2]}
    }
    \and
    \inferrule{
      \T[1] \simeq \T*[1] \\
      \T[2] \simeq \T*[2]
    }{
      \T[1] \TPair \T[2]
      \simeq
      \T*[1] \TPair \T*[2]
    }
    \and
    \inferrule{
      \T[1] \simeq \T*[1] \\
      \T[2] \simeq \T*[2]
    }{
      \TSum{\T[1]}{\T[2]}
      \simeq
      \TSum{\T*[1]}{\T*[2]}
    }
  \end{mathpar}
  \caption{Modified rules for $\simeq$}
  \label{fig:cfstb-equivalence}
\end{figure}

% \begin{definition}[Substituting definitions]
%   Let $F$ be a FreeST session type and $U$ be a mapping from variables to FreeST
%   session types such that $\text{freevars}(F) \subseteq \text{dom}(U)$. We write
%   $F[U]$ for the FreeST type obtained by substituting the free variables in $F$
%   according to mapping $U$.
%   A well formed substitution has no free variables. 
% \end{definition}

Let $F$ be a FreeST session type and $U$ be a mapping from variables to FreeST
session types. In the following, we write $F[U]$ for the pair of $F$ and $U$
where $\text{freevars}(F) \subseteq \text{dom}(U)$. This notation includes
unfolding of definitions in the following way:
\begin{gather*}
  \frac{(X \mapsto F) \in U}{X[U] \simeq F[U]}
\end{gather*}

\begin{lemma}[Soundness of the translation]
  If $\T[1] \simeq \T[2]$, then $F_1 [U_1] \simeq F_2 [U_2]$
  where $\langle F_1, U_1 \rangle = \tofreest{\T[1]}$ and
  $\langle F_2, U_2 \rangle = \tofreest{\T[2]}$.
  \label{lem:sound-translation}
\end{lemma}
\begin{proof}
  We sketch a proof. The soundness proof requires two things: the translation of
  annotations in the arrow and product types preserves equivalence and the
  equivalence of the explicit recursion with the $\mu$-binder syntax. For the
  former we can perform induction on the equivalence of $\T[1] \simeq \T[2]$ and
  use equivalence and bisimulation lemmas from \citet{DBLP:conf/icfp/ThiemannV16}
  for regular type and session type cases, respectively. The pairs we introduce
  for annotations can be handled trivially by the type equivalence rules. To
  show explicit recursion translation preserves equivalence, we can use the laws
  for $\mu$-types from \citet{DBLP:conf/icfp/ThiemannV16}. 
\end{proof}

\begin{lemma}[Completeness of the translation]
  Let $\langle F_1, U_1 \rangle = \tofreest{\T[1]}$ and
  $\langle F_2, U_2 \rangle = \tofreest{\T[2]}$. If
  $F_1 [U_1] \simeq F_2 [U_2]$, then $\T[1] \simeq \T[2]$.
\end{lemma}
\begin{proof}
  Similar to \ref{lem:sound-translation} but the opposite direction.
\end{proof}

\clearpage

\section{Metatheory}
\label{sec:appendix-metatheory}

\begin{figure}[h]
  \begin{mathpar}
    \RuleExpValueBlocked \and
    \RuleExpConstBlocked \and
    % \RuleParReady \and
    % \RuleNuReady \and
    \RuleParBlocked \and
    \RuleNuBlocked \and
    \RuleNuBlockedAcqLeft \and
    \RuleNuBlockedAcqRight \and
    \RuleNuBlockedAcqBoth
  \end{mathpar}
  where
  \[
    \begin{array}{ll}
      \BlockedChannelsSend & \BlockedChannelsRecv \\
      \BlockedChannelsSelect & \BlockedChannelsBranch \\
      \BlockedChannelsTerm & \BlockedChannelsWait \\
      \BlockedChannelsAcquire & \BlockedChannelsDrop \\
      \BlockedChannelsLSplit & \BlockedChannelsRSplit \\
      \BlockedChannelsFork & \BlockedChannelsNew \\
      \BlockedChannelsDiscard \\\\
      \AcqBlockedChannelsSend & \AcqBlockedChannelsRecv \\
      \AcqBlockedChannelsSelect & \AcqBlockedChannelsBranch \\
      \AcqBlockedChannelsTerm & \AcqBlockedChannelsWait \\
      \AcqBlockedChannelsAcquire & \AcqBlockedChannelsDrop \\
      \AcqBlockedChannelsLSplit & \AcqBlockedChannelsRSplit \\
      \AcqBlockedChannelsFork & \AcqBlockedChannelsNew \\
      \AcqBlockedChannelsDiscard 
    \end{array}
  \]
  \caption{Blocked Processes (Blocked $\Proc$)}
  \label{fig:blocked-expr-proc}
\end{figure}


%%% Local Variables:
%%% mode: latex
%%% TeX-master: "../popl2027-supplement"
%%% End:
